\documentclass[final]{beamer}
\usetheme{Warsaw}
\usepackage[utf8]{inputenc}
\usepackage[english,russian]{babel}
\usepackage{mathpartir,stmaryrd}
\usepackage[
  orientation=portrait,size=a1,scale=1.25
]{beamerposter}
\geometry{hmargin=2.5cm}
\linespread{1.15}
\usepackage{poster/beamerthemesharelatex}

\title[]{\bf \huge
Погружение уточняющих типов данных
в системы с зависимыми типами}
\author{Соколов Павел Павлович}
\institute{ФКН НИУ ВШЭ, Москва}
\date{27 октября 2025}

\begin{document}
\begin{frame}[t]
\begin{multicols}{2}

\section{Основные понятия}

В \textit{теории типов} основным объектом изучения
являются {\color{color4}системы типов} --- формальные
математические модели языков программирования (либо
чисто логических исчислений), состоящие из описания
грамматических форм языка, правил вычисления, и,
наконец, {\color{color4}правил типизации},
классифицирующих выражения по тому, какие ``значения''
они вычисляют (понятие значения, конечно, крайне
условно и зависит от языка). Особый интерес представляют
системы типов, позволяющие одновременно кодировать и
алгоритмы, и доказательство их корректности. Здесь
помогает принятие парадигмы {\color{color4} Утверждения
есть (некоторые) типы}, отождествляющее понятия из левой
колонки в таблице ниже с понятиями из правой колонки:

\vskip1ex
\begin{table}
\centering
\caption{Парадигма Propositions as (Some) Types}
\begin{tabular}{c|c}
Логика & Теория типов\\
\hline
Утверждение $P$ & Тип $P$ \\
Теорема $T$ & Населённый тип $T$ \\
Предположение $H$ & Переменная $x:H$ \\
Доказательство $\Gamma\vdash T$
& Программа $\Gamma\vdash \pi:T$ \\
Устранение сечений & Нормализация вычислений \\
Непротиворечивость
& $\nexists t.\;\cdot\vdash t:\bot$ \\
$\vdash P \land Q$
& $\vdash (p : P, q : Q) : P \times Q$ \\
$\vdash P \to Q$
& $\vdash(\lambda (x:P).q:Q):P\to Q$ \\
Кванторы & Зависимые типы \\
$\vdash \exists x \in X. P(x)$
& $\vdash (w : X, \pi : P(w)) : \Sigma_{x:X}P(x)$ \\
$\vdash \forall x \in X. P(x)$
& $\vdash(\lambda(x:X).\pi:P(x)):\Pi_{x:X}P(x)$ \\
\end{tabular}
\end{table}
\vskip1ex

В рамках этой парадигмы поиск доказательства сводится к
синтезу программы подходящего типа; корректность
доказательства проверит компилятор. Правда, выписывать
такое доказательство трудно и долго, так что желательно
максимально автоматизировать и упростить этот процесс.

\section{Историческая справка}

Первой успешной практической реализацией языка на основе
парадигмы Propositions as Types следует считать систему
Rocq (ранее Coq) Тьерри Коканда, основанную на
Исчислении Конструкций~\cite{COQUAND198895}. В этой
системе постулируется существование двух сортов:
$Prop$ (сорт утверждений) и $Type$ (сорт типов данных),
типов (зависимых) функций, связывающих сорта, а также
\underline{индуктивных} определений (с их помощью можно
как постулировать новые утверждения, так и строить
пользовательские типы данных). Другие аналогичные
системы~\cite{norell:thesis, moura2021lean} основываны
на той же самой базовой идее, но проблема автоматизации
доказательств решается немного по-разному.

\begin{itemize}
    \item В Rocq~\cite{the_coq_development_team_2024_14542673}
        автоматизация достигается посредством
        \underline{тактик}, специальных команд,
        генерирующих собственно код доказательств;
    \item В Lean~\cite{moura2021lean} тоже используют
        тактики, а также охотно пользуются аксиомами
        классической логики и теории множеств, что
        приближает доказательства в Lean к обычной,
        ``бумажной'' математике;
    \item В Agda~\cite{norell:thesis} тактик нет;
        вместо этого эксплуатируется рефлексия и
        продвинутый алгоритм \underline{элаборации},
        позволяющий программисту опускать большинство
        технических аннотаций типов.
\end{itemize}

В частности, автоматизация в Agda больше
сконцентрирована на реализации \underline{солверов} ---
алгоритмов, генерирующих доказательство в рамках
некоторой ограниченной теории (ведь для всей системы
типов целиком поиск доказательства --- неразрешимая
задача, как и для классической логики первого порядка).
С наглядным примером реализации этого подхода можно
ознакомиться в~\cite{agdareflection}.

Альтернативный подход --- изначально ограничить систему
типов таким образом, чтобы в ней были выразимы только
разрешимые предикаты; кроме этого, желательно объединить
программу и доказываемые про неё свойства в единую
сущность: $\{ x : T \;|\; \varphi(x) \}$. Такие типы
называются \underline{уточняющими}; гибкость по
сравнению с обычными зависимыми типами достигается
посредством следующего
правила~\cite{shen2021smtbasedrefinementtypesagda}:
\begin{mathpar}
\inferrule{
    \Gamma\vdash t:\{x:T\;|\;\varphi(x)\}\\
    \llbracket\Gamma,x:T\rrbracket\models
    \varphi(x)\Rightarrow\psi(x)
}{\Gamma\vdash t:\{x:T\;|\;\psi(x)\}}
\end{mathpar}
Обязательство по доказательству $\varphi\Rightarrow\psi$
возлагается на компилятор; в статье, на которую мы
ссылаемся выше, предлагается использовать любой готовый
SMT-солвер. К сожалению, такое решение не расширяемо:
набор автоматически доказуемых утверждений ограничен
возможностями используемого солвера. Кроме этого,
полученная система не самодостаточна.

\section{Основные результаты}

Мы предлагаем систему с зависимыми типами,
поддерживающую уточняющие типы данных на уровне
библиотеки, где множество автоматически доказуемых
утверждений расширяемо пользователем. При этом система
не опирается на внешний источник доказательств, а строит
их самостоятельно. Это достигается с помощью алгоритма
\underline{разрешения ограничений}~\cite{selsam2020tabledtypeclassresolution},
используемого в компиляторах функциональных языков
программирования для поддержки перегрузки операторов.

Из интересных деталей, по сравнению с системой, кратко
описанной в~\cite{shen2021smtbasedrefinementtypesagda},
мы добавляем следующее правило:
\begin{mathpar}
\inferrule{
    \Gamma\vdash t:\{x:T\;|\;\phi(x)\}
    \\\Gamma\vdash u:\{y:U(t:T)\;|\;\psi(t, y)\}
}{\Gamma\vdash (t,u)
    :\{(x,y):\Sigma_{x:T}U(x)
    \;|\;\phi(x)\land\psi(x, y)\}}
\end{mathpar}
Благодаря ему, система получает возможность рассуждать
о совокупностях объектов (например, об алгебраических
структурах и операциях на них) эргономичным для
пользователя образом.

\subsection{Идеи доказательств}

Доказательство сохранения типизации производится
стандартным образом.

Логическая непротиворечивость следует из сводимости
нашей системы к Исчислению
конструкций~\cite{COQUAND198895}, про которое была
доказана~\cite{setsintypes} взаимная консистентность с
ZFC (+ необходимое количество недостижимых кардиналов).

Ключевой деталью в построении сведения является
разрешимость процедуры поиска доказательств, благодаря
которой сведение получается эффективным.

\section{Выводы и дальнейшие направления исследований}

Таким образом, мы получили исчисление, в рамках которого
автоматизируется поиск тривиальных заключений из
доказанных пользователем теорем. Интересно также
рассмотреть следующие расширения:

\vskip1ex
\begin{itemize}
    \item Практическую реализацию системы;
    \item Интегрирование полноценных солверов
        разрешимых теорий внутрь системы;
    \item Взаимодействие уточняющих типов с
        гомотопическими теориями типов;
    \item Оптимизация сложности генерируемых
        доказательств;
    \item Взаимодействие уточняющих типов с эффективной
        компиляцией.
\end{itemize}
\vskip1ex

\subsection{Список литературы}
\bibliographystyle{ieeetr}
\bibliography{bibliography}

\end{multicols}
\end{frame}
\end{document}
